% SPDX-License-Identifier: CC-BY-4.0
\section{Evidence and the research framework}\label{sec:framework}

\subsection{Development and checks}

The results were developed in a public research notebook~\cite{Notebook}. Its
dated entries record the exact computations, the failed routes and the reviews.
\begin{itemize}
 \item \textbf{Exact computation.} Before the general statement was proved,
 minimum degrees were computed exactly for \(n\le7\) by rank computations
 over \(\Ft\). This produced the formula.
 \item \textbf{Earlier upper bound.} Until the literature search found
 Menezes' construction, the upper bound was proved only for \(k\le31\), by
 explicit symmetric factors found by a solver and checked independently.
 \item \textbf{Checks of the construction.} The construction of
 Theorem~\ref{thm:per-order} was checked by exact multiplicity computation in
 417 cases with \(k\le12\) and \(n\le7\), 305 of them below Menezes' range
 \(n\ge k-1\).
 \item \textbf{Review.} The informal proofs of the step, the lower bound and the
 construction were checked by fresh-context reviewers before being recorded.
 \item \textbf{Formalization.} The whole chain of Theorem~\ref{thm:main} was
 then formalized. The Lean proofs follow the arguments given here.
\end{itemize}
The one-dimension step, the truncated product and the Schwartz--Zippel lemma
use generalized forms, such as every \(n\ge1\), every \(s\) and every field; this
paper states them in those forms.

\subsection{Noemesis}

The work used \emph{Noemesis}, an open framework for autonomous mathematical
research, distributed with the repository~\cite{Repo}. It consists of
instructions, tools and record-keeping conventions that let an AI assistant
continue an investigation across sessions with inspectable evidence:
\begin{itemize}
 \item a living overview and a claim index;
 \item dated, append-only research entries, including failures;
 \item protected and timed computations with retained outputs;
 \item dependency-indexed Lean formalization, linked from the claim index.
\end{itemize}
These mechanisms make the work recoverable and reviewable. They do not make a
recorded proof correct; the Lean formalization does that for the statements
carrying [Lean] links.

\subsection{Contributions and AI assistance}\label{sec:model-assistance}

\begin{itemize}
 \item \textbf{The earlier note~\cite{Note}.} Proposition~\ref{prop:cover-bound}
 was drafted with GPT-6 Astra in the Codex agent.
 \item \textbf{This paper.} The rest of the mathematics was produced by Claude
 Opus 5.5 in Claude Code, working inside the framework above under the
 author's direction: the exact computations, the formula, the one-dimension
 step, the literature search that found~\cite{Men26}, the characteristic-two
 proof of the construction, the complete Lean formalization, and the draft of
 this paper. Fresh-context Claude Opus reviewers checked the informal
 arguments.
 \item \textbf{The author.} The author chose the problem, directed the
 investigation and designed the research framework.
\end{itemize}
The paper has not been externally peer reviewed. Formal verification covers
the statements with [Lean] links, under the definitions in the linked files.
It does not cover the prose or the comparison with the literature.

The construction of Section~\ref{sec:upper} is due to Menezes~\cite{Men26}, and
the multiplicity Schwartz--Zippel lemma to Dvir, Kopparty, Saraf and
Sudan~\cite{DKSS13}. Original manuscript material is distributed under the
\href{https://creativecommons.org/licenses/by/4.0/}{Creative Commons Attribution
4.0 International license (CC BY 4.0)}.
